\section{Discussion}

The results do not support naive hard-to-predict word masking as a reliable improvement over random masking. HHM is mechanically viable and respects the fairness constraint, but it does not improve validation loss in the completed three-seed comparison. Raw-loss CMS-Morph also fails to improve validation loss, even though it assigns a larger adaptive share than Accuracy-Morph. The negative result is therefore not simply that the adaptive share was too small; the difficulty signal itself matters.

The official BabyLM diagnostics make the result more nuanced. HHM is slightly higher than the random baseline on each available official zero-shot metric, even though its validation loss is worse. These gains are small and not statistically established, but they show that hard-to-predict masking can alter measured linguistic behavior without improving the local MLM objective. This is a useful diagnostic outcome: validation loss and downstream-style BabyLM diagnostics need not move in lockstep under small masking-policy changes.

The strongest lesson is that raw high loss is a noisy online difficulty signal. A high-loss token may be rare, ambiguous, poorly tokenized, or unlearnable in context. Reallocating masks toward such tokens does not necessarily create better pretraining signal. Smoothed prediction correctness appears more stable in the current experiments because it estimates repeated prediction failure over exposures rather than reacting directly to loss magnitude. Accuracy-Morph is therefore the most promising tested signal: across seeds 1 and 2 it improves matched validation loss by $-0.0155 \pm 0.0051$ and Entity Tracking by $+5.51 \pm 3.27$. This is encouraging pilot evidence, not a confirmed positive result.

The entity-augmented follow-up is also instructive. Adding repeated-token pressure, an intermediate difficulty band, and late adaptive decay does not preserve the Accuracy-Morph advantage. This suggests that future work should not simply add more heuristics. The most defensible next step is a multi-seed Accuracy-Morph evaluation with cleaner prediction-correctness scoring, modest token and morph adaptive shares, and carefully tested difficulty-band filtering.

Relative to the closest ACL Anthology work, the main difference is experimental scope. Prior adaptive-masking papers show that changing masked content, schedules, task-guided selection, or predictability-aware token probabilities can improve pretraining or BabyLM performance \citep{gu2020train,lee2022efficient,yang2023learning,wilf2023difference,edman2025mask}. BabyLM system papers show that larger changes, such as hybrid causal-masked objectives, lightweight architectures, curated data, optimizer changes, or masked diffusion, can be stronger overall routes to leaderboard performance \citep{hu2024babylm,charpentier2024gptbert,yu2024antlm,haller2025sample,kosmopoulou2025masked}. Our experiments are deliberately narrower. The value of the paper is not that HHM beats those systems; it does not. The value is that it isolates the masking-policy variable and documents which online signals survive that constraint. Under this stricter test, repeated-token and raw-loss signals are insufficient, while correctness-smoothed token and trigram sketches are a plausible low-cost direction.
